\documentclass[12pt,a4paper]{ctexart}

% ============ 页面设置 ============
\usepackage[margin=2.5cm]{geometry}

% ============ 数学公式 ============
\usepackage{amsmath,amssymb}

% ============ 图片支持 ============
\usepackage{graphicx}
\usepackage{float}

% ============ 表格 ============
\usepackage{booktabs}
\usepackage{array}

% ============ 超链接 ============
\usepackage[hidelinks]{hyperref}

% ============ 代码高亮 ============
\usepackage{listings}
\usepackage{xcolor}

\lstset{
    basicstyle=\small\ttfamily,
    breaklines=true,
    frame=single,
    numbers=left,
    numberstyle=\tiny,
    backgroundcolor=\color{gray!5}
}

% ============ 标题信息 ============
\title{\heiti 需求分析文档 \\[6pt] \large —— XX系统}
\author{\songti 小组成员：张三、李四、王五}
\date{\today}

\begin{document}

\maketitle
\thispagestyle{empty}
\newpage

% ============ 目录 ============
\tableofcontents
\newpage

% ============ 正文 ============

\section{引言}
\subsection{编写目的}
本文档旨在明确XX系统的功能需求和非功能需求，为后续的设计、开发和测试阶段提供依据。

\textbf{预期读者}包括：
\begin{itemize}
    \item 项目开发团队
    \item 项目管理人员
    \item 系统最终用户
    \item 测试人员
\end{itemize}

\subsection{项目背景}
简要描述项目产生的背景、要解决的问题以及项目的目标和范围。

\subsection{术语定义}
\begin{table}[H]
\centering
\caption{术语与缩写说明}
\begin{tabular}{lll}
\toprule
\textbf{术语/缩写} & \textbf{全称} & \textbf{说明} \\
\midrule
SRS & Software Requirements Specification & 软件需求规格说明 \\
UML & Unified Modeling Language & 统一建模语言 \\
CRUD & Create, Read, Update, Delete & 增删改查操作 \\
\bottomrule
\end{tabular}
\end{table}

\section{系统概述}
\subsection{系统目标}
用1-2段话描述系统要达成的核心目标。

\subsection{用户角色}
\begin{table}[H]
\centering
\caption{用户角色定义}
\begin{tabular}{lll}
\toprule
\textbf{角色} & \textbf{职责} & \textbf{权限级别} \\
\midrule
管理员   & 系统配置、用户管理、数据维护     & 最高 \\
普通用户 & 核心业务操作                   & 中等 \\
访客     & 浏览公开信息                   & 最低 \\
\bottomrule
\end{tabular}
\end{table}

\subsection{运行环境}
\begin{itemize}
    \item \textbf{操作系统}：Windows 10/11、Linux、macOS
    \item \textbf{数据库}：MySQL 8.0 / PostgreSQL
    \item \textbf{后端框架}：待定
    \item \textbf{前端框架}：待定
\end{itemize}

\section{功能需求}
\subsection{功能模块总览}
\begin{figure}[H]
\centering
（系统功能模块图待补充）
\caption{系统功能模块图}
\end{figure}

\subsection{模块一：用户管理}
\subsubsection{用户注册}
\begin{itemize}
    \item \textbf{编号}：FUNC-USER-001
    \item \textbf{优先级}：高
    \item \textbf{描述}：新用户通过填写用户名、邮箱、密码完成注册
    \item \textbf{输入}：用户名（3-20字符）、邮箱（合法格式）、密码（8位以上）
    \item \textbf{输出}：注册成功 → 跳转登录页；注册失败 → 显示错误提示
    \item \textbf{前置条件}：用户名和邮箱未被占用
    \item \textbf{后置条件}：用户记录写入数据库，发送验证邮件
\end{itemize}

\subsubsection{用户登录}
\begin{itemize}
    \item \textbf{编号}：FUNC-USER-002
    \item \textbf{优先级}：高
    \item \textbf{描述}：已注册用户通过用户名/邮箱 + 密码登录系统
    \item \textbf{输入}：用户名或邮箱、密码
    \item \textbf{输出}：登录成功 → 进入主页；登录失败 → 提示错误
\end{itemize}

\subsection{模块二：核心业务（待补充）}
根据项目实际需求在此处补充核心业务功能模块。

\section{非功能需求}
\subsection{性能需求}
\begin{itemize}
    \item 页面响应时间不超过\underline{2秒}
    \item 支持\underline{100}名用户同时在线
    \item 数据库查询响应时间不超过\underline{500ms}
\end{itemize}

\subsection{安全性需求}
\begin{itemize}
    \item 密码使用bcrypt等算法加密存储
    \item 敏感数据传输使用HTTPS
    \item 防止SQL注入、XSS、CSRF攻击
    \item 用户操作日志记录
\end{itemize}

\subsection{可用性需求}
\begin{itemize}
    \item 系统7×24小时可用，全年可用性 $\ge$ 99\%
    \item 界面简洁直观，符合用户操作习惯
    \item 提供清晰的错误提示和帮助文档
\end{itemize}

\subsection{可维护性需求}
\begin{itemize}
    \item 代码遵循统一编码规范
    \item 关键模块有单元测试覆盖
    \item 提供完整的API文档
\end{itemize}

\section{用例描述}
\subsection{用例图}
（用例图待补充，可使用TikZ-UML或外部工具绘制）

\subsection{用例规约示例}
\begin{table}[H]
\centering
\caption{用例UC-001：用户注册}
\begin{tabular}{lp{11cm}}
\toprule
\textbf{用例编号} & UC-001 \\
\textbf{用例名称} & 用户注册 \\
\textbf{参与者}   & 未注册用户 \\
\textbf{前置条件} & 用户进入注册页面 \\
\textbf{后置条件} & 用户账号创建成功 \\
\midrule
\multicolumn{2}{l}{\textbf{基本事件流：}} \\
\multicolumn{2}{l}{\parbox{11cm}{
1. 用户点击"注册"按钮 \\
2. 系统显示注册表单 \\
3. 用户填写用户名、邮箱、密码 \\
4. 用户点击"提交" \\
5. 系统验证输入合法性 \\
6. 系统创建用户账号 \\
7. 系统发送验证邮件 \\
8. 系统提示注册成功 \\
}} \\
\midrule
\multicolumn{2}{l}{\textbf{异常事件流：}} \\
\multicolumn{2}{l}{\parbox{11cm}{
3a. 用户名已被占用 → 提示更换用户名 \\
3b. 邮箱格式不合法 → 提示正确格式 \\
3c. 密码强度不足 → 提示密码要求 \\
}} \\
\bottomrule
\end{tabular}
\end{table}

\section{数据需求}
\subsection{数据实体}
列出系统涉及的核心数据实体及其关系，例如：
\begin{itemize}
    \item \textbf{用户}：UserID, Username, Email, Password, Role, CreatedAt
    \item \textbf{其他实体}：根据项目补充
\end{itemize}

\subsection{ER图}
（ER图待补充）

\section{验收标准}
\begin{enumerate}
    \item 所有高优先级功能通过测试
    \item 性能指标达到非功能需求要求
    \item 安全测试无高危漏洞
    \item 用户文档和API文档完整
    \item 源代码通过代码审查
\end{enumerate}

\section{附录}
\subsection{参考文献}
\begin{enumerate}
    \item Ian Sommerville, \textit{Software Engineering}, 10th Edition, Pearson, 2015.
    \item 相关课程资料
\end{enumerate}

\subsection{修订历史}
\begin{table}[H]
\centering
\caption{文档修订历史}
\begin{tabular}{llll}
\toprule
\textbf{版本} & \textbf{日期} & \textbf{修订人} & \textbf{说明} \\
\midrule
v0.1 & \today & 小组 & 初稿 \\
\bottomrule
\end{tabular}
\end{table}

\end{document}
